\documentclass[10pt,twocolumn]{article}
\usepackage[margin=0.9in]{geometry}
\usepackage{amsmath,amssymb,mathtools}
\usepackage{booktabs}
\usepackage{microtype}
\usepackage{hyperref}

\title{A Conservative Galilean Frame-Tracking Strategy for Long-Duration TRML Simulations}
\author{AthenaK TRML Collaboration}
\date{\today}

\begin{document}
\maketitle

\begin{abstract}
Long-duration turbulent radiative mixing-layer (TRML) calculations are often limited by slow, cumulative interface drift rather than by local numerical instability. We present a conservative Galilean frame-tracking strategy that stabilizes the thermally selected interface without introducing non-conservative source forcing. The controller combines filtered global measurements with cadence-gated proportional--derivative feedback, optional integral bias cancellation, and explicit command limiting. A two-stage supercomputer campaign (96 production runs) was used to calibrate operating parameters across $\xi=100$ and $\xi=1000$. The resulting configuration improves centering and robustness relative to baseline while preserving practical computational cost. We report the recommended parameter set, summarize the empirical ranking criteria, and document the tradeoff between aggressive centering and conservative boundary behavior.
\end{abstract}

\section{Introduction}
In TRML calculations, physically meaningful long-time interpretation depends on maintaining the mixing interface within the statistically representative region of the domain. Even when local hydrodynamic updates remain stable, small systematic biases can produce secular vertical migration over hundreds of dynamical times. This drift contaminates integrated diagnostics, shortens useful averaging windows, and can trigger unphysical interactions with vertical boundaries.

The objective of frame tracking is therefore not to ``force'' the flow, but to maintain a computational reference frame that follows the evolving layer. The method must satisfy three requirements simultaneously: (i) act on a physically coherent global signal, (ii) preserve conservative finite-volume updates, and (iii) remain robust in strongly turbulent, intermittently multi-phase states.

\section{Control Formulation}
The tracker is implemented as a discrete-time Galilean velocity translation applied at controlled cadence. At each eligible control event, the algorithm selects cells in a temperature band,
\begin{equation}
T_{\mathrm{lo}} \le T \le T_{\mathrm{hi}},
\end{equation}
and forms density-weighted global moments for vertical position and velocity. The position observable may be the centroid, the thermal-band midpoint, or their convex blend. To suppress high-frequency turbulence-induced jitter, both observables are evolved with exponential moving averages.

The command law is a bounded PD controller with optional integral leakage term:
\begin{equation}
v_{\mathrm{cmd}} = \mathcal{L}\!\left(v_P + v_D + v_I\right),
\end{equation}
where $\mathcal{L}$ denotes the ordered limiter sequence (absolute cap, optional predicted-position clamp, slew-rate cap, and final absolute cap). The applied quantity is a uniform Galilean shift $\Delta v_z$.

Crucially, the hydrodynamic state update is conservative: momentum is translated exactly and, for ideal EOS runs, total energy is updated with the corresponding Galilean cross and quadratic terms. The tracker therefore preserves the conservative structure of the underlying solver while controlling drift of the thermally defined layer.

\section{Campaign Design and Calibration Strategy}
Parameter calibration was performed using two campaigns on Frontier:
\begin{enumerate}
  \item \textbf{Broad sweep} (2026-03-02): 48 runs spanning baseline-adjacent and stress-test settings.
  \item \textbf{Follow-up refinement} (2026-03-03): 48 runs concentrated near the robust region identified by campaign 1.
\end{enumerate}
Both campaigns used $\chi=10^{1.75}$, $t_{\mathrm{lim}}=150$, and paired runs at $\xi=100$ and $\xi=1000$.

Ranking emphasized centering quality over the post-transient window ($t\ge20$), then penalized pathological controller behavior:
\begin{equation}
Q = |\langle z_{\mathrm{ft}}\rangle| + w_s f_{\mathrm{skip}} + w_c f_{\mathrm{cap}} + w_\sigma\,\sigma(z_{\mathrm{ft}}).
\end{equation}
To avoid selecting numerically aggressive configurations, we added a conservative diagnostic check based on interface-edge contact timing from history diagnostics.

\section{Results}
The refinement campaign substantially improved reliability and controller behavior relative to the broad sweep. The first campaign produced 43 successful windows out of 48 cases, whereas the second produced 48/48. Mean skip fraction decreased from 0.1168 to 0.0187, and mean cap fraction from 0.0941 to 0.0121.

The lowest raw quality score in the follow-up set corresponded to a more aggressive parameter combination. However, conservative diagnostics showed earlier edge contact, indicating reduced safety margin for long integrations. The selected production point therefore favored a slightly less aggressive but more conservative operating regime with near-identical runtime cost.

Relative to the follow-up baseline, the selected point reduced paired quality-max by 16.0\% and paired quality-mean by 5.2\%, while improving the conservative edge-contact timing metric by 52.4\%.

\section{Recommended Operating Point}
Table~\ref{tab:best} reports the recommended production configuration.

\begin{table}[t]
\centering
\small
\caption{Recommended frame-tracking parameters for production TRML runs.}
\label{tab:best}
\begin{tabular}{@{}ll@{}}
\toprule
Parameter & Value \\
\midrule
\texttt{ft\_apply\_every} & 2 \\
\texttt{ft\_mode} & pd \\
\texttt{ft\_tau\_avg} & 0.04 \\
\texttt{ft\_tau\_relax} & 2.6 \\
\texttt{ft\_tau\_vel} & 0.09 \\
\texttt{ft\_temp\_lo} & 0.050 \\
\texttt{ft\_temp\_hi} & 0.066 \\
\texttt{ft\_max\_abs\_boost} & 0.012 \\
\texttt{ft\_max\_boost\_change} & 0.002 \\
\texttt{ft\_tau\_int} & 12.0 \\
\texttt{ft\_int\_max\_abs} & 0.0015 \\
\texttt{ft\_int\_leak\_tau} & 8.0 \\
\texttt{ft\_int\_unsat\_only} & true \\
\bottomrule
\end{tabular}
\end{table}

\section{Cadence-Cost Tradeoff}
A frequent practical question is whether much larger actuation cadence can be used to reduce cost. The campaign data do not support this for production quality. Moving from \texttt{ft\_apply\_every}=2 to 3--4 reduced measured runtime only modestly, but degraded ranking quality metrics strongly (about 34--76\%, depending on statistic and $\xi$). This pattern indicates that very sparse actuation is a false economy for long-time statistical studies.

\section{Conclusions}
A conservative Galilean frame-tracking controller can materially improve long-duration TRML interpretability without compromising conservative finite-volume updates. The calibrated operating point reported here is not the most aggressive centering configuration; it is the most defensible production configuration once robustness, conservative boundary behavior, and compute cost are considered jointly. This distinction is essential for publication-grade simulation campaigns in which statistical fidelity over long integration windows is the scientific objective.

\end{document}
